\subsection{CPython Architectural Evolution}

Program-visible \texttt{dict} behavior---keyed lookup, insertion-order iteration since 3.7---rests on a compact hash-table layout introduced for memory efficiency (PEP 468 implementation track; insertion-order preservation became a language guarantee in 3.7).

\subsubsection{The Classic Dictionary Model: Sparse Hash Table Storage and Historically Unspecified Iteration Order}

Pre-compact \texttt{dict} objects stored entries directly in a sparse \texttt{PyDictEntry[]} array: each slot carried hash, key pointer, and value pointer (24 bytes per slot on 64-bit builds). Most slots stayed empty to preserve open-addressing probe termination. Iteration order was an implementation artifact of slot occupancy, not a semantic promise.

\subsubsection{The Modern Compact Dictionary: Dense Entry Storage, Sparse Indexing, Memory Reduction, and Insertion-Order Preservation}

The compact layout splits responsibilities:

\begin{itemize}
    \item \textbf{Sparse index array} --- a byte- or integer-index table sized to the next power of two; each cell holds an index into the dense entry block, or an empty/dummy sentinel. Probing walks this lightweight array.
    \item \textbf{Dense \texttt{PyDictKeyEntry[]} block} --- contiguous records \texttt{(hash, key)} stored in insertion order; values live in a parallel array indexed by the same dense offset.
\end{itemize}

Lookup hashes the key, probes the index array, follows the index into the dense block, compares hash integers, then calls \texttt{\_\_eq\_\_}. Iteration walks the dense block sequentially, yielding insertion order as a structural side effect without maintaining a separate linked list. Compared with the monolithic sparse entry table, this representation typically saves on the order of 20--40\% of dict heap footprint for small-to-medium mappings.

Updating an existing key replaces its value in place without moving the dense index; deleting and re-inserting a key appends a new dense slot at the end, moving that key to the tail of iteration order.

\subsubsection{Algorithmic Efficiency Matrix: Amortized \texttt{O(1)} Complexity for Key Insertion, Value Retrieval, and Structural Deletion}

Insert, indexed lookup, and delete by key each perform one hash computation plus an expected constant number of index probes and occasional dense-block appends. Table growth triggers a rehash of all live entries into a larger index array---the same amortized pattern as sets. Dicts remain the right structure when the access pattern is keyed associative fetch, not positional pair scanning.
